%% ============================================================
%% AULA 4 — Treinando a rede: risco empírico e descida de gradiente
%% GBC073 — Inteligência Computacional
%% Baseada na Lecture 4 de 11-785 Introduction to Deep Learning (CMU),
%% traduzida, condensada e modernizada (ênfase em PyTorch).
%%
%% Compilar:  latexmk -xelatex aula04.tex   (XeLaTeX, duas passadas)
%% ============================================================
\documentclass[aspectratio=169, 11pt]{beamer}

\makeatletter
\def\input@path{{./}{../}}
\makeatother

\usetheme{InteligenciaComputacional}   % ANTES do babel: fontspec vem primeiro
\usepackage{babel}
\babelprovide[import, main]{brazilian}

% ---- código Python/PyTorch -------------------------------------------------
\usepackage{listings}
\lstdefinestyle{python}{
  language=Python,
  basicstyle=\ttfamily\scriptsize,
  keywordstyle=\color{icAzul}\bfseries,
  commentstyle=\color{icCinza}\itshape,
  stringstyle=\color{icVerde},
  emph={torch,nn,optim,backward,step,zero_grad},
  emphstyle=\color{icMagenta},
  showstringspaces=false,
  columns=fullflexible,
  keepspaces=true,
  frame=leftline,
  framerule=2pt,
  rulecolor=\color{icAzul!40},
  xleftmargin=1em,
}
\lstset{style=python}

\title[Aula 4 — Treinando a rede]{Aula 4: Treinando a rede — risco empírico e descida de gradiente}
\subtitle[GBC073 — Inteligência Computacional]{Otimização \textbullet\ Gradientes \textbullet\ Funções de perda \textbullet\ PyTorch}
\author[Prof.\ Marcelo Albertini]{Prof.\ Marcelo Keese Albertini}
\institute{GBC073 — Inteligência Computacional \textbullet\ Universidade Federal de Uberlândia}
\date{2026}

\begin{document}

% ------------------------------------------------------------ capa
\begin{frame}[plain, noframenumbering]
  \titlepage
\end{frame}

% ------------------------------------------------------------ roteiro
\begin{frame}{Onde estamos}
  \framesubtitle{Da aula passada para esta}

  \begin{ideiabox}
    Já sabemos que uma rede neural \emph{pode representar} praticamente qualquer
    função (aproximação universal). A pergunta desta aula é outra:
    \alert{como encontrar os pesos} que fazem a rede representar a função certa?
  \end{ideiabox}

  \medskip
  Plano da aula:
  \begin{itemize}\setlength{\itemsep}{0.2ex}
    \item Treinar $=$ minimizar o \textbf{risco empírico} — um problema de otimização
    \item Ferramentas: derivadas, gradientes, Hessiana
    \item O algoritmo: \textbf{descida de gradiente} (à mão e com \texttt{autograd})
    \item Montar o problema para redes: a rede $f$, as saídas, a \textbf{divergência}
    \item O resultado que prepara a retropropagação: $\nabla_{\vet{z}}\,\mathrm{div} = \hat{\vy} - \vy$
  \end{itemize}
\end{frame}

% ============================================================
\section{Do risco empírico à otimização}

\begin{frame}{Recapitulando: o problema de aprender}
  \begin{defbox}{Minimização do risco empírico (ERM)}
    Dado um conjunto de treino $\D = \{(\vx_1, \vy_1), \dots, (\vx_T, \vy_T)\}$,
    escolha os parâmetros que minimizam a divergência média entre a saída da
    rede e o alvo:
    \[
      \widehat{\vw} = \argmin_{\vw}\; \erro(\vw), \qquad
      \erro(\vw) = \frac{1}{T} \sum_{i=1}^{T}
        \mathrm{div}\big( f(\vx_i;\, \vw),\; \vy_i \big)
    \]
  \end{defbox}

  A \textbf{perda} $\erro(\vw)$ é \emph{empírica}: não conhecemos a
  distribuição real dos dados, só as amostras de $\D$.

  \begin{ideiabox}
    Treinar uma rede \emph{é} minimizar uma função. Todo o deep learning moderno
    se apoia em uma única ideia: se $f$ e $\mathrm{div}$ forem
    \alert{diferenciáveis}, dá para minimizar $\erro(\vw)$ seguindo o gradiente.
  \end{ideiabox}
\end{frame}

% ============================================================
\section{Derivadas, gradientes e mínimos}

\begin{frame}{Derivada: o multiplicador local}
  \begin{defbox}{Derivada como aproximação linear}
    Em resolução fina o bastante, toda função suave é localmente linear:
    \[
      f(x + \Delta x) \;\approx\; f(x) + f'(x)\, \Delta x
    \]
    A derivada diz \emph{quanto a saída muda} por unidade de mudança minúscula na
    entrada.
  \end{defbox}

  \begin{itemize}\setlength{\itemsep}{0.2ex}
    \item Sinal: $f' > 0$ $\Rightarrow$ aumentar $x$ aumenta $f$;\;
          $f' < 0$ $\Rightarrow$ diminui;\; $f' = 0$ $\Rightarrow$ localmente plana.
          Magnitude $|f'|$: inclinação da curva
  \end{itemize}

  \begin{ideiabox}
    É exatamente esta leitura que o treino usa: o gradiente de $\erro$ em relação
    a um peso diz se \emph{aumentar aquele peso} aumenta ou diminui o erro — e o
    quanto.
  \end{ideiabox}
\end{frame}

\begin{frame}{Gradiente: derivada de função escalar de vetor}
  \begin{defbox}{Gradiente}
    Para $f : \R^n \to \R$:\;
    $\nabla_{\vx} f = \big[ \deriv{f}{x_1} \;\cdots\; \deriv{f}{x_n} \big]^{\!\top}$,
    \; com \; $\Delta f \approx \nabla_{\vx} f \cdot \Delta\vx$
  \end{defbox}

  \begin{columns}[T]
    \begin{column}{0.52\textwidth}
      \begin{teobox}{Direção de crescimento mais rápido}
        O produto interno $\nabla f \cdot \Delta\vx$ é máximo quando $\Delta\vx$
        está \emph{alinhado} com $\nabla f$. Logo, o gradiente aponta para onde
        $f$ \alert{cresce mais rápido} — e $-\nabla f$ para onde decresce mais
        rápido.
      \end{teobox}
    \end{column}
    \begin{column}{0.44\textwidth}
      \begin{center}
        \begin{tikzpicture}[scale=0.72]
          \draw[icAzul!30, thick] (0,0) ellipse (2.6 and 1.25);
          \draw[icAzul!55, thick] (0,0) ellipse (1.7 and 0.82);
          \draw[icAzul!85, thick] (0,0) ellipse (0.85 and 0.41);
          \fill[icAzul] (0,0) circle (1.6pt);
          \node[below, font=\scriptsize, text=icAzul] at (0,-0.05) {mínimo};
          \node[circle, fill=icMagenta, inner sep=1.6pt] (p) at (1.7,0.82) {};
          \draw[sinapse destac] (p) -- ++(0.95,0.46)
            node[above, font=\scriptsize, text=icMagenta] {$\nabla f$};
          \draw[sinapse] (p) -- ++(-0.95,-0.46)
            node[below, font=\scriptsize, text=icCinza] {$-\nabla f$};
        \end{tikzpicture}

        {\scriptsize\color{icCinza} curvas de nível: $\nabla f$ é perpendicular a elas}
      \end{center}
    \end{column}
  \end{columns}
\end{frame}

\begin{frame}{Mínimos: pontos críticos e a Hessiana}
  \begin{teobox}{Condições de mínimo local (função suave)}
    Em um mínimo local de $f : \R^n \to \R$ valem duas condições:
    \(\nabla f = \vet{0}\) (ponto crítico) e a \textbf{Hessiana}
    \(\nabla^2 f\), matriz das segundas derivadas
    \(\big[\nabla^2 f\big]_{ij} = \partial^2 f / \partial x_i \partial x_j\),
    é \alert{positiva semidefinida} (autovalores $\geq 0$).
  \end{teobox}

  \begin{itemize}\setlength{\itemsep}{0.2ex}
    \item Gradiente nulo sozinho não basta: máximos, mínimos, selas e inflexões
          \emph{todos} têm $\nabla f = \vet{0}$ (em 1D: $f''>0$ em mínimos,
          $f''<0$ em máximos)
  \end{itemize}

  \begin{alertabox}
    Em redes profundas ninguém verifica a Hessiana (milhões$^2$ de entradas):
    aceita-se o ponto onde a descida de gradiente estaciona. Em alta dimensão,
    pontos de \alert{sela} são muito mais comuns que mínimos ruins.
  \end{alertabox}
\end{frame}

\begin{frame}{Pergunta para conduzir em aula}
  \begin{quizbox}
    $f(\vx)$ é uma função escalar de um vetor $\vx$. Partindo de $\vx_0$, em qual
    direção devemos dar um passo para obter a \emph{maior redução} de $f$?

    \smallskip
    \opcoes{Na direção de $\nabla f$ \;/\; Perpendicular a $\nabla f$ \;/\; Oposta a $\nabla f$}
  \end{quizbox}

  \pause

  \begin{respbox}
    \textbf{Oposta ao gradiente.} O gradiente é a direção de subida mais rápida;
    seu oposto é a descida mais rápida. Perpendicular ao gradiente o valor de $f$
    não muda (é a direção da curva de nível). Este é o coração do algoritmo da
    próxima seção.
  \end{respbox}
\end{frame}

% ============================================================
\section{Descida de gradiente}

\begin{frame}{Por que soluções iterativas?}
  Para minimizar, poderíamos resolver $\nabla_{\vw} \erro(\vw) = \vet{0}$
  analiticamente\dots

  \begin{itemize}\setlength{\itemsep}{0.2ex}
    \item Funciona para funções simples (regressão linear tem solução fechada);
          para uma rede com milhões de pesos e não linearidades compostas, o
          sistema $\nabla \erro = \vet{0}$ é \alert{intratável}
  \end{itemize}

  \begin{ideiabox}
    A saída: começar de um \emph{chute} $\vw^{(0)}$ e refiná-lo passo a passo,
    sempre na direção que reduz a perda. Duas decisões definem o método:
    \emph{em que direção} andar e \emph{qual o tamanho} do passo.
  \end{ideiabox}

  \begin{alertabox}
    Todo o treinamento moderno — do Perceptron ao Transformer — é variante desta
    ideia; o que muda é como o passo é calculado (SGD, momento, Adam\dots).
  \end{alertabox}
\end{frame}

\begin{frame}{O algoritmo}
  \begin{algbox}{Descida de gradiente (\emph{gradient descent})}
    \begin{itemize}\setlength{\itemsep}{0.2ex}
      \item Inicialize $\vw^{(0)}$; $k \leftarrow 0$
      \item Enquanto não convergir:
            \(\;\vw^{(k+1)} \leftarrow \vw^{(k)} - \taxa\, \nabla_{\vw}\, \erro\big(\vw^{(k)}\big)\)
      \item Critérios de parada:
            \(\big|\erro(\vw^{(k+1)}) - \erro(\vw^{(k)})\big| < \epsilon_1\)
            \;ou\; \(\norm{\nabla_{\vw} \erro}^2 < \epsilon_2\)
    \end{itemize}
  \end{algbox}

  Anatomia do passo:
  \[
    \vw \leftarrow \vw
      \;-\; \termo{\taxa}{taxa de aprendizado}
      \;\cdot\; \termo{\nabla_{\vw}\, \erro(\vw)}{gradiente da perda}
  \]

  \begin{itemize}\setlength{\itemsep}{0.2ex}
    \item $\taxa$ pequena demais: converge devagar \quad\textbullet\quad
          $\taxa$ grande demais: oscila ou diverge
    \item Para funções \textbf{convexas} (formato de tigela) e $\taxa$ adequada,
          converge ao mínimo \alert{global}; caso contrário, a um mínimo local ou sela
  \end{itemize}
\end{frame}

\begin{frame}[fragile]{Em código: descida de gradiente à mão}
  Minimizar $f(x) = x^2 - 4x + 5$, cujo mínimo é $x^* = 2$
  (derivada $f'(x) = 2x - 4$):

\begin{lstlisting}
import torch

x   = torch.tensor(8.0)        # chute inicial
eta = 0.1                      # taxa de aprendizado
for t in range(60):
    grad = 2*x - 4             # derivada calculada a mao
    x    = x - eta * grad      # passo: oposto ao gradiente
print(x)                       # tensor(2.0000)
\end{lstlisting}

  \begin{ideiabox}
    Cada linha do laço é a fórmula do slide anterior. O único trabalho
    ``inteligente'' foi derivar $f$ à mão — e é exatamente isso que o
    \texttt{autograd} elimina.
  \end{ideiabox}
\end{frame}

\begin{frame}[fragile]{Em código: o mesmo, com autograd}
\begin{lstlisting}
x   = torch.tensor(8.0, requires_grad=True)
opt = torch.optim.SGD([x], lr=0.1)

for t in range(60):
    loss = x**2 - 4*x + 5      # constroi o grafo de computacao
    opt.zero_grad()            # zera gradientes acumulados
    loss.backward()            # autograd: calcula d(loss)/dx
    opt.step()                 # x <- x - lr * x.grad
\end{lstlisting}

  \begin{itemize}\setlength{\itemsep}{0.2ex}
    \item \texttt{loss.backward()} percorre o grafo de trás para frente e preenche
          \texttt{x.grad} — é a \alert{retropropagação}, tema da próxima aula;
          \texttt{opt.step()} aplica $\vw \leftarrow \vw - \taxa \nabla \erro$
  \end{itemize}

  \begin{alertabox}
    Esquecer \texttt{opt.zero\_grad()} é o bug clássico: o PyTorch \emph{acumula}
    gradientes entre chamadas de \texttt{backward()}, e o passo passa a usar a
    soma de todos os gradientes já calculados.
  \end{alertabox}
\end{frame}

\begin{frame}{Convexidade: a garantia e a realidade}
  \begin{columns}[T]
    \begin{column}{0.48\textwidth}
      \begin{tcolorbox}[iccaixa=icCiano, title={O que a teoria garante}]
        \begin{itemize}\setlength{\itemsep}{0.2ex}
          \item Função convexa $+$ passo adequado $\Rightarrow$ mínimo global
          \item Base sólida: regressão linear, logística
        \end{itemize}
      \end{tcolorbox}
    \end{column}
    \begin{column}{0.48\textwidth}
      \begin{tcolorbox}[iccaixa=icMagenta, title={O que o deep learning vive}]
        \begin{itemize}\setlength{\itemsep}{0.2ex}
          \item Perda de rede profunda: \alert{não convexa}
          \item Milhões de mínimos locais e selas
        \end{itemize}
      \end{tcolorbox}
    \end{column}
  \end{columns}

  \medskip
  \begin{ideiabox}
    Por que funciona mesmo assim? Em redes \emph{superparametrizadas}, quase
    todos os mínimos locais têm perda próxima da ótima. O problema prático não é
    ``cair no mínimo errado'', e sim caminhar de forma eficiente na paisagem —
    daí os otimizadores modernos.
  \end{ideiabox}
\end{frame}

% ============================================================
\section{Montando o problema para redes neurais}

\begin{frame}{Três coisas a definir}
  \[
    \widehat{\vw} = \argmin_{\vw}\; \frac{1}{T} \sum_{i=1}^{T}
      \mathrm{div}\big( f(\vx_i;\, \vw),\; \vy_i \big)
  \]

  \begin{columns}[T]
    \begin{column}{0.31\textwidth}
      \begin{tcolorbox}[iccaixa=icAzul, title={A rede $f(\cdot\,; \vw)$}]
        \small Qual arquitetura? Quais são os parâmetros $\vw$?
      \end{tcolorbox}
    \end{column}
    \begin{column}{0.31\textwidth}
      \begin{tcolorbox}[iccaixa=icCiano, title={Os pares $(\vx_i, \vy_i)$}]
        \small Como representar entradas e saídas como números?
      \end{tcolorbox}
    \end{column}
    \begin{column}{0.31\textwidth}
      \begin{tcolorbox}[iccaixa=icMagenta, title={A divergência}]
        \small Como medir o erro entre $\hat{\vy}$ e $\vy$? Precisa ser diferenciável!
      \end{tcolorbox}
    \end{column}
  \end{columns}

  \medskip
  O resto da aula responde às três perguntas, uma por vez.
\end{frame}

\begin{frame}{A rede $f$: o MLP e sua notação}
  \begin{columns}[T]
    \begin{column}{0.54\textwidth}
      \begin{itemize}\setlength{\itemsep}{0.3ex}
        \item Rede \emph{direta} (\emph{feedforward}) em camadas, sem laços;
              a entrada é a camada $0$ (sem neurônios: são os dados)
        \item Saída do neurônio $j$ da camada $k$: $y_j^{(k)}$;\;
              $\pesoij{i}{j}^{(k)}$: peso do neurônio $i$ da camada $k{-}1$ ao
              neurônio $j$ da camada $k$; viés $b_j^{(k)}$
        \item Cada neurônio: $\campo = \vw^{\!\top}\vx + b$,\; $y = \ativ(\campo)$
              — afim $+$ ativação \alert{diferenciável}
      \end{itemize}
    \end{column}
    \begin{column}{0.42\textwidth}
      \begin{center}
        \scalebox{0.52}{\mlp{4,5,5,3}}

        {\color{icCinza}\scriptsize entrada em ciano, ocultas em azul, saída em rosa}
      \end{center}
    \end{column}
  \end{columns}

  \begin{ideiabox}
    Todos os pesos e vieses juntos formam o vetor gigante $\vw$ do gradiente.
    Rede moderna: de milhões a trilhões de coordenadas, mesmo algoritmo.
  \end{ideiabox}
\end{frame}

\begin{frame}{Ativações e suas derivadas}
  A descida de gradiente exige derivar tudo — inclusive as ativações:
  \begin{center}\small
    $\sig(z) = 1/(1+e^{-z})$,\; $\sig' = \sig(1-\sig)$
    \qquad
    $\tanh' = 1 - \tanh^2$
    \qquad
    $\relu(z) = \max(0, z)$,\; $\relu'(z) = [z > 0]$
  \end{center}

  \begin{columns}[T]
    \begin{column}{0.48\textwidth}
      \begin{tcolorbox}[iccaixa=icCinza, title={Histórico}]
        \small Sigmoide e $\tanh$ nas ocultas dominaram até $\sim$2012.
        Saturam: derivada $\approx 0$ longe da origem.
      \end{tcolorbox}
    \end{column}
    \begin{column}{0.48\textwidth}
      \begin{tcolorbox}[iccaixa=icVerde, title={Moderno}]
        \small Família ReLU (ReLU, GELU, SiLU) é o padrão nas ocultas:
        não satura no lado positivo.
      \end{tcolorbox}
    \end{column}
  \end{columns}

  \begin{alertabox}
    Ativação saturada $\Rightarrow$ gradiente $\approx 0$ $\Rightarrow$ o peso
    \alert{para de aprender} (\emph{desvanecimento do gradiente}) — foi o que
    tirou a sigmoide das ocultas; hoje ela vive quase só na saída de
    classificadores binários.
  \end{alertabox}
\end{frame}

\begin{frame}{Ativações vetoriais: softmax}
  \begin{defbox}{Softmax}
    Ativação que acopla todas as saídas da camada, transformando valores afins
    $z_1, \dots, z_K$ (os \emph{logits}) em uma distribuição de probabilidade:
    \[
      \hat{y}_i = \frac{e^{z_i}}{\sum_{j=1}^{K} e^{z_j}}
      \qquad\Longrightarrow\qquad
      \hat{y}_i > 0, \;\; \sum_i \hat{y}_i = 1
    \]
  \end{defbox}

  \begin{itemize}\setlength{\itemsep}{0.2ex}
    \item Mudar \emph{um} peso afeta \emph{todas} as saídas (por isso ``vetorial'');
          uso padrão na saída de classificadores multiclasse:
          $\hat{y}_i = $ probabilidade da classe $i$
  \end{itemize}

  \begin{ideiabox}
    O softmax reaparece no coração dos Transformers: a \emph{atenção} usa esta
    mesma função para converter escores de similaridade em pesos que somam 1.
  \end{ideiabox}
\end{frame}

\begin{frame}{Representando as saídas $\vy$}
  \begin{itemize}\setlength{\itemsep}{0.4ex}
    \item \textbf{Regressão} (saída real): use o valor diretamente; um neurônio
          por dimensão, saída linear
    \item \textbf{Binária} (``é um gato?''): alvo $y \in \{0, 1\}$, saída
          sigmoide $\hat{y} \in (0,1)$ lida como $P(\text{classe} = 1)$
    \item \textbf{Multiclasse}: alvo em representação \emph{one-hot}, saída softmax
  \end{itemize}

  \begin{defbox}{Vetor one-hot}
    Para $K$ classes em ordem fixa, a classe $c$ vira o vetor com 1 na posição
    $c$ e 0 nas demais. Com classes [gato, cão, camelo, chapéu, flor]:
    \[
      \text{camelo} \;\mapsto\; \vy = (0,\, 0,\, 1,\, 0,\, 0)^{\!\top}
    \]
  \end{defbox}

  \begin{itemize}\setlength{\itemsep}{0.2ex}
    \item A rede produz $\hat{\vy}$: vetor de $K$ probabilidades que somam 1
    \item O ideal seria $\hat{\vy}$ igual ao one-hot; na prática, quanto mais massa
          na classe certa, melhor
  \end{itemize}
\end{frame}

\begin{frame}[fragile]{Em código: a rede $f$ em PyTorch}
  Classificador de dígitos (entrada $28 \times 28 = 784$ pixels, 10 classes):

\begin{lstlisting}
import torch.nn as nn

modelo = nn.Sequential(
    nn.Linear(784, 256), nn.ReLU(),   # camada oculta 1
    nn.Linear(256, 128), nn.ReLU(),   # camada oculta 2
    nn.Linear(128, 10),               # saida: 10 logits
)

logits = modelo(x)                    # x: tensor (B, 784)
\end{lstlisting}

  \begin{itemize}\setlength{\itemsep}{0.2ex}
    \item Cada \texttt{nn.Linear} guarda uma matriz de pesos e um viés — as
          coordenadas de $\vw$; \texttt{modelo.parameters()} entrega todas ao otimizador
    \item Repare: \alert{não há softmax na saída}. A rede devolve \emph{logits} —
          o porquê fica claro em dois slides
  \end{itemize}
\end{frame}

% ============================================================
\section{Divergências: medindo o erro}

\begin{frame}{Regressão: divergência L2}
  \begin{defbox}{Divergência L2 (erro quadrático)}
    Para saídas reais, a escolha padrão é a distância euclidiana ao quadrado
    (escalada):
    \[
      \mathrm{div}(\hat{\vy}, \vy) = \frac{1}{2} \norm{\vy - \hat{\vy}}^2
             = \frac{1}{2} \sum_i (y_i - \hat{y}_i)^2
      \qquad\Longrightarrow\qquad
      \deriv{\mathrm{div}}{\hat{\vy}} = \hat{\vy} - \vy
    \]
  \end{defbox}

  \begin{itemize}\setlength{\itemsep}{0.2ex}
    \item Diferenciável, e a derivada é o \alert{erro} — o $\tfrac{1}{2}$ só
          cancela o 2 da derivada. Em PyTorch: \texttt{nn.MSELoss()}
  \end{itemize}

  \begin{alertabox}
    Para \emph{classificação}, a L2 é ruim: com saída sigmoide/softmax gera
    gradientes minúsculos justo quando a rede está muito errada. Daí a próxima
    divergência.
  \end{alertabox}
\end{frame}

\begin{frame}{Classificação: KL e entropia cruzada}
  \begin{defbox}{Entropia cruzada (\emph{cross-entropy})}
    Para alvo one-hot $\vy$ (classe correta $c$) e saída softmax $\hat{\vy}$:
    \[
      \mathrm{div}(\hat{\vy}, \vy)
        = -\sum_i y_i \log \hat{y}_i
        = -\log \hat{y}_c
      \qquad
      \deriv{\mathrm{div}}{\hat{y}_i} =
        \begin{cases} -1/\hat{y}_c & i = c \\ 0 & i \neq c \end{cases}
    \]
  \end{defbox}

  \begin{itemize}\setlength{\itemsep}{0.2ex}
    \item Minimizar $-\log \hat{y}_c$ $=$ maximizar a probabilidade da classe
          correta (máxima verossimilhança); para alvo one-hot, KL e entropia
          cruzada \alert{coincidem}
  \end{itemize}

  \begin{ideiabox}
    A entropia cruzada \emph{explode} quando $\hat{y}_c \to 0$: gradiente enorme
    justo quando a rede está confiante e errada. Esse castigo acelera a
    convergência — é a perda de todo modelo de classificação e linguagem moderno.
  \end{ideiabox}
\end{frame}

\begin{frame}{Pergunta para conduzir em aula}
  \begin{quizbox}
    A entropia cruzada atinge seu mínimo quando $\hat{\vy} = \vy$. Portanto,
    nesse ponto, a derivada $\partial\,\mathrm{div} / \partial \hat{y}_c$ é zero
    — como acontece com a L2.

    \smallskip
    \opcoes{Verdadeiro \;/\; Falso}
  \end{quizbox}

  \pause

  \begin{respbox}
    \textbf{Falso.} No mínimo ($\hat{y}_c = 1$) a derivada vale $-1/\hat{y}_c = -1
    \neq 0$. O mínimo existe porque $\hat{y}_c$ não pode passar de 1 (restrição do
    softmax), não porque a curva aplaina. Consequência prática: o gradiente
    continua ``empurrando'' $\hat{y}_c$ para cima mesmo perto do acerto — em
    relação aos \emph{logits}, porém, o gradiente $\hat{\vy} - \vy$ zera direitinho.
  \end{respbox}
\end{frame}

\begin{frame}{O resultado central da aula}
  \framesubtitle{O gradiente em relação aos logits é o erro}

  Encadeando softmax $+$ entropia cruzada e derivando em relação aos valores
  afins $\vet{z}$ (os logits) que entram no softmax, quase tudo cancela:
  \[
    \mathrm{div} = -\log \hat{y}_c, \quad \hat{\vy} = \mathrm{softmax}(\vet{z})
    \qquad\Longrightarrow\qquad
    \mdestaque{icMagenta}{\nabla_{\vet{z}}\, \mathrm{div} = \hat{\vy} - \vy}
  \]

  \begin{itemize}\setlength{\itemsep}{0.2ex}
    \item O mesmo vale para saída linear $+$ L2: gradiente $= \hat{\vy} - \vy$
    \item Interpretação limpa: o que se propaga para trás a partir da saída é
          simplesmente \alert{o erro} — probabilidade prevista menos alvo
  \end{itemize}

  \begin{ideiabox}
    Não é coincidência: L2$\,\leftrightarrow\,$saída linear e entropia
    cruzada$\,\leftrightarrow\,$softmax são pares ``casados''. É o primeiro elo
    da cadeia que a retropropagação estende por toda a rede na próxima aula.
  \end{ideiabox}
\end{frame}

\begin{frame}[fragile]{Em código: a perda em PyTorch}
\begin{lstlisting}
criterio = nn.CrossEntropyLoss()  # log-softmax + NLL, tudo em um

logits = modelo(x)                # (B, 10) -- SEM softmax antes!
perda  = criterio(logits, alvo)   # alvo: indices (B,), nao one-hot

perda.backward()                  # gradiente flui: comeca em y_hat - y
\end{lstlisting}

  \begin{alertabox}
    Erro clássico: aplicar \texttt{nn.Softmax} na rede \emph{e} usar
    \texttt{CrossEntropyLoss}. A perda já embute o softmax (em forma
    numericamente estável, via \texttt{log-sum-exp}); softmax duplo achata os
    gradientes e o treino ``anda de lado''.
  \end{alertabox}

  \begin{itemize}\setlength{\itemsep}{0.2ex}
    \item \texttt{CrossEntropyLoss(label\_smoothing=0.1)}: troca o one-hot por
          alvo suavizado ($1-\epsilon$ na classe certa, $\epsilon/(K{-}1)$ nas demais) —
          truque moderno de regularização, usado em Transformers
  \end{itemize}
\end{frame}

% ============================================================
\section{Fechamento}

\begin{frame}{História até aqui}
  \begin{itemize}\setlength{\itemsep}{0.4ex}
    \item Redes neurais são aproximadores universais; treiná-las é
          \textbf{minimizar o risco empírico} — um problema de otimização
    \item O \textbf{gradiente} aponta a subida mais rápida; a
          \textbf{descida de gradiente} caminha no sentido oposto:
          $\vw \leftarrow \vw - \taxa \nabla_{\vw} \erro$
    \item Definimos a rede ($f$ diferenciável), as saídas (real / sigmoide /
          one-hot $+$ softmax) e a divergência (\textbf{L2}; \textbf{entropia
          cruzada}); nos pares casados, o gradiente na saída é o erro:
          $\nabla_{\vet{z}}\, \mathrm{div} = \hat{\vy} - \vy$
  \end{itemize}

  \begin{exbox}{1 — Softmax e entropia cruzada}
    Com logits $\vet{z} = (2,\, 0,\, -1)$ e classe correta $c = 1$:
    (a) calcule $\hat{\vy} = \mathrm{softmax}(\vet{z})$;
    (b) calcule a perda $-\log \hat{y}_1$;
    (c) calcule $\nabla_{\vet{z}}\,\mathrm{div} = \hat{\vy} - \vy$ e verifique
    que as componentes somam zero. Confira com \texttt{torch}.
  \end{exbox}

  \medskip
  \centering {\color{icCinza}\small Próxima aula: \textbf{retropropagação} — como
  levar $\hat{\vy} - \vy$ da saída até cada peso da rede.}
\end{frame}

\end{document}
